\label{sec:docker-summary}

Im Rahmen dieses Kapitels wurde die Containerisierung und das Deployment der
Anwendung umfassend dargelegt. Ziel war es, einen praxisorientierten und
transparenten Einblick in die Nutzung zeitgemäßer Container-Technologien zu
vermitteln und deren Nutzen im Kontext eines realen Schulprojekts zu
demonstrieren. Der Schwerpunkt lag dabei gezielt auf Docker, Docker Compose sowie
dem Deployment auf einer virtuellen Maschine.

Zunächst wurden die Fundamentals der Containerisierung erläutert. Dabei wurde
ersichtlich, dass Container eine leichtgewichtige, portable Alternative zu
klassischen Virtualisierungskonzepten repräsentieren. Durch die Kapselung von
Anwendung und Abhängigkeiten ermöglichen sie reproduzierbare Builds und
konsistente Ausführungsumgebungen, unabhängig vom eingesetzten Host-System.

Im Anschluss wurde der konkrete Docker-Einsatz im Projekt dargestellt. Sowohl
für das Backend als auch für das Frontend wurden dedizierte Dockerfiles erstellt,
die auf dem Prinzip des Multi-Stage Builds basieren. Diese Methodik ermöglicht
eine klare Separation von Build- und Laufzeitumgebung, wodurch die finalen Images
kompakt gehalten und Sicherheitsrisiken minimiert werden. Die Integration der
Dockerfiles als Code-Listings steigert zudem die Nachvollziehbarkeit der
Implementierung.

Darauf aufbauend wurde Docker Compose als zentrales Instrument für das Deployment
mehrerer Services präsentiert. Durch die deklarative Beschreibung der
Anwendungsarchitektur lassen sich Frontend, Backend, Datenbank und weitere
Hilfsdienste gemeinsam initialisieren und administrieren. Insbesondere die
Spezifikation von Abhängigkeiten, Netzwerken und Volumes trägt zu einem stabilen
und wartungsfreundlichen Betrieb bei. Das Compose-Setup konstituiert damit die
technische Basis für den Produktivbetrieb auf der virtuellen Maschine.

Das Deployment auf der virtuellen Maschine stellt einen wesentlichen Schritt dar,
um die Anwendung realitätsnah zu betreiben. Die Kombination aus Virtualisierung
und Containerisierung vereint die Vorteile beider Ansätze: einerseits eine
klar abgegrenzte Serverumgebung, andererseits flexible und leicht wartbare
Anwendungskomponenten. Persistente Volumes gewährleisten, dass essenzielle Daten
auch bei Neustarts oder Aktualisierungen erhalten bleiben.

Abschließend wurde Kubernetes als alternative Orchestrierungslösung eingeordnet
und bewusst vom Projekt abgegrenzt. Obwohl Kubernetes für umfangreiche, skalierbare
und hochverfügbare Systeme zahlreiche Vorteile bietet, würde dessen Einsatz den
Rahmen und die Anforderungen des vorliegenden Projekts signifikant überschreiten.
Die Entscheidung für Docker Compose repräsentiert daher eine bewusste Abwägung
zwischen Komplexität und Nutzen.

\subsubsection{Ausblick}

Für eine potenzielle Weiterentwicklung des Projekts ergeben sich mehrere
Ansatzpunkte. Eine naheliegende Erweiterung wäre die Integration eines
automatisierten CI/CD-Prozesses, bei dem Docker-Images nach jeder Modifikation
automatisch generiert, getestet und in eine Container Registry transferiert werden.
Dadurch ließe sich der Deployment-Prozess weiter optimieren und Fehlerquellen
minimieren.

Darüber hinaus könnte bei steigender Nutzerzahl oder expandierender
Anwendungsarchitektur eine Migration zu Kubernetes erwogen werden. Die bereits
implementierte Containerisierung und eindeutige Separation der Services bilden
dafür eine solide Basis. Kubernetes würde in diesem Szenario zusätzliche
Möglichkeiten zur Skalierung, Ausfallsicherheit und Lastverteilung eröffnen.

Zusammenfassend demonstriert das Projekt, dass mit Docker und Docker Compose
bereits ohne komplexe Orchestrierungslösungen ein stabiler, wartungsfreundlicher
und zeitgemäßer Produktivbetrieb realisiert werden kann. Die gewählte Architektur
eignet sich sowohl für den aktuellen Projektumfang als auch für künftige
Erweiterungen und stellt damit eine nachhaltige Lösung dar.